\subsection*{NT-R08. A discriminant-uniform lower logarithmic density}

Retain the finite etale setting of NT-R07 and assume only RH for
$\zeta_L$. There is an absolute constant $C>0$ such that
\[
 V_X\le C r^2\log^2(\Delta_L+2),\qquad
 \underline\delta\{\Pi_X(x)<r\operatorname{Li}_2(x)\}
 \ge \frac1{1+C\log^2(\Delta_L+2)}.
 \tag{NT-R08a}
\]
Here $\underline\delta$ denotes the lower logarithmic density
defined by the left side of (NT-R07a). The same constant works
for every number field $K$, every finite Galois splitting field
$L/K$, and every nonempty finite etale cover with this splitting
field. No lower bound on zero spacing or upper bound on individual
multiplicities independent of $\Delta_L$ is assumed. The
almost-periodic limiting process is still taken separately for
each fixed extension; no uniform rate of convergence to the law
is asserted.

\paragraph{Complete proof.}
Write $n=[L:\mathbb Q]$ and $D=\Delta_L$. Let
$m_L(\gamma)=\operatorname{ord}_{s=1/2+i\gamma}\zeta_L(s)$.
For every $\gamma>0$, the cyclic-character formula (NT-R05a)
and nonnegativity of the permutation character $t$ give
\[
 |a_t(1/2+i\gamma)|
 \le\sum_C\frac{|C|}{|G|}t(C)
       \sum_{\chi\in\widehat H_C}
          \operatorname{ord}_{1/2+i\gamma}L(s,\chi,L/F_C)
 =r\,m_L(\gamma).
 \tag{NT-R08b}
\]
This is an exact inequality: the inner sum equals the order
of $\zeta_L$, and the outer sum equals $r$ by Burnside.

The unconditional zero-count formula in Grenie--Molteni,
AppendixA TheoremA.1, author p.26, specialized to the trivial
character of $L$, implies with an absolute implied constant
\[
 U_j:=\sum_{\substack{\gamma\ {\rm distinct}\\j<\gamma\le j+1}}
                  m_L(\gamma)
       \ll \log D+n\log(j+3)\qquad(j=0,1,2,\ldots).
 \tag{NT-R08c}
\]
For clarity, the source counts both signs and includes
multiplicities. Its main term is
$T\{\log D+n\log(T/(2\pi e))\}/\pi$, with error
$O(\log D+n\log T+n)$ for $T\ge1$, and an additional
$O(n)$ constant term. For $j\ge2$ subtract its values at
$j-1$ and $j+1$; the difference of the main terms has the
displayed size by differentiation, and the sum of the two
error bounds has that size as well. This bounds the positive
ordinates in $(j,j+1]$. For $j=0,1$ use the source bound for
the total count at $T=3$. The absolute constants in its table
therefore give (NT-R08c) uniformly in the field.

Within a unit interval the orders are nonnegative, so the sum
of their squares is at most $U_j^2$. Equations (NT-R06a) and
(NT-R08b) yield
\[
 \begin{split}
 V_X
 &\le 2r^2\sum_{\gamma>0}
                  \frac{m_L(\gamma)^2}{1/4+\gamma^2}\\
 &\ll r^2\sum_{j=0}^{\infty}
          \frac{(\log D+n\log(j+3))^2}{1/4+j^2}\\
 &\ll r^2(\log D+n)^2.
 \end{split}
 \tag{NT-R08d}
\]
The final constant is absolute, since the series with numerator
$1$, $\log(j+3)$ or $\log^2(j+3)$ all converge against this
denominator.

We include the elementary removal of the degree parameter.
Minkowski's discriminant bound gives
\[
 \sqrt D\ge\frac{n^n}{n!}\left(\frac{\pi}{4}\right)^{r_2}.
\]
An exact reference is Milne, \emph{Algebraic Number Theory},
v3.08, Theorem4.3 printed p.70 and the proof of Theorem4.9
printed p.72. For $n\ge2$, the arithmetic-geometric mean
inequality gives $n!\le((n+1)/2)^n$, and $r_2\le n/2$.
Consequently
\[
 \sqrt D\ge
 \left(\frac{2n}{n+1}\sqrt{\frac{\pi}{4}}\right)^n
 \ge\left(\frac{2\sqrt\pi}{3}\right)^n.
\]
Since $2\sqrt\pi/3>1$, this proves
$n\ll\log D$ with an absolute constant for $n\ge2$.
The case $n=1,D=1$ is handled by $\log(D+2)$.
Substitution in (NT-R08d) proves the first assertion of
(NT-R08a). Finally $B_X\ge r$ in NT-R07 gives
\[
 \frac{B_X^2}{B_X^2+V_X}
 \ge \frac1{1+C\log^2(D+2)},
\]
which proves the lower-density assertion.

The variance estimate is consistent with the squared
$\|t\|_1^2$ bound in Fiorilli--Jouve Proposition4.6(87),
author p.40 and its proof p.41. The proof above derives the
coarse bound directly from the uniform zero count and gives
the finite etale consequence without invoking a bounded
multiplicity hypothesis. No numerical value for the absolute
constant is asserted.

In the Grothendieck language of NT-F02 and NT-R07, this is a
uniform lower-frequency statement for the Frobenius trace on
the permutation realization of $H^0$ of a finite etale cover.
The generic rank of that realization may grow; the final lower
bound depends on the splitting-field discriminant, not that
rank. It does not give a uniform threshold in $x$ or a uniform
error term for the approach to the limiting law.
